\section{Introduction}

Stability conditions organize the objects of a derived category by a complex
central charge and a real phase, with finite filtrations into semistable
objects. Their physical antecedent is Douglas's proposal of $\Pi$-stability
for D-branes \cite{Douglas2001}; Bridgeland gave the mathematical framework
used here \cite{Bridgeland2007}. At a Gepner point, one seeks a stability
condition compatible with a distinguished autoequivalence: the functor should
rotate the central charge and translate every semistable phase by a fixed
amount. We construct such a condition for every smooth quintic threefold.

Let $X\subset\mathbb P^4_{\C}$ be a smooth quintic hypersurface, and put
$\D=\Db(\Coh(X))$ and $H=c_1(\OO_X(1))$. The structure sheaf is spherical:
$\Ext^*(\OO_X,\OO_X)=\C\oplus\C[-3]$. Using the spherical twist of
Seidel--Thomas \cite{SeidelThomas2001}, we consider the autoequivalence
\[
 \Phi=T_{\OO_X}\circ(-\otimes\OO_X(1)),\qquad
 T_{\OO_X}(E)=\operatorname{Cone}\bigl(
 R\Hom(\OO_X,E)\otimes\OO_X\longrightarrow E\bigr).
\]
We write $\Knum(X)$ for the Grothendieck group modulo the radical of its Euler
pairing, and $\Knum(X)_{\R}=\Knum(X)\otimes_{\Z}\R$.
A numerical central charge is a homomorphism $Z\colon\Knum(X)\to\C$.

Following Bridgeland \cite{Bridgeland2007}, a slicing $\mathcal P$ is a family
of full additive subcategories $\mathcal P(\varphi)\subset\D$, indexed by
$\varphi\in\R$, with
$\mathcal P(\varphi+1)=\mathcal P(\varphi)[1]$,
vanishing $\Hom(\mathcal P(\varphi_1),\mathcal P(\varphi_2))$ for
$\varphi_1>\varphi_2$, and finite Harder--Narasimhan filtrations by distinguished
triangles with strictly decreasing phases. A nonzero object of
$\mathcal P(\varphi)$ has charge in $\R_{>0}e^{i\pi\varphi}$. Local finiteness
requires the extension categories of sufficiently short phase intervals to be
quasi-abelian of finite length. We use the support property of
Kontsevich--Soibelman \cite[Section~1.2]{KontsevichSoibelman2008}, expressed on
the full numerical group: there exist a norm on $\Knum(X)_{\R}$ and a
constant $C>0$ such that
\[
 \norm{[E]}\le C|Z(E)|
 \quad\text{for every nonzero }E\in\mathcal P(\varphi).
\]
The pair $(Z,\mathcal P)$ is a numerical Bridgeland stability condition when
these charge and slicing requirements hold, with local finiteness.
We will also verify directly that each phase category is closed under direct
summands, using the phase cuts constructed in Lemma~\ref{lem:phase-cuts}.

Toda formulated Gepner-type stability for graded matrix factorizations
\cite[Definition~2.3 and Conjecture~1.2]{Toda2013}, and gave the normalized
quintic formulation in \cite[Conjecture~3.1]{TodaQuintic2013}.
The following theorem positively resolves this quintic conjecture.

\begin{theorem}\label{thm:main}
Every smooth complex quintic threefold admits a numerical Bridgeland stability
condition $\sigma=(Z,\mathcal P)$ with the support property on $\Knum(X)_{\R}$
such that, for every $E\in\D$ and $\varphi\in\R$,
\[
 Z(\Phi E)=e^{2\pi i/5}Z(E),\qquad
 \Phi\bigl(\mathcal P(\varphi)\bigr)=\mathcal P(\varphi+2/5),\qquad
 Z(\OO_x)=-1
\]
for every closed point $x\in X$.
\end{theorem}

The normalization identifies our central charge with Toda's by the uniqueness
of the normalized eigenform proved in Proposition~\ref{prop:central-charge}.

\paragraph{Prior work and the remaining construction.}
The functor relation $\Phi^5\simeq[2]$ and its eigencharge are established
features of the quintic Gepner problem. Orlov's equivalence
\cite[Theorem~2.5]{Orlov2009} and the compatibility with grade shift proved by
Ballard--Favero--Katzarkov \cite[Proposition~5.8 and Remark~5.12]{BallardFaveroKatzarkov2012}
identify the relevant categorical symmetry; Toda recalls this identification
in \cite[Theorem~2.5]{TodaQuintic2013}. A direct computation with
Fourier--Mukai kernels appears in Aspinwall \cite[Section~7.1.4]{Aspinwall2004}.
Section~\ref{sec:numerical} fixes the numerical conventions used in our
construction; Appendix~\ref{app:periodicity} supplies a relative-kernel proof
of the functor relation.

The double-tilt approach to threefold stability was developed by
Bayer--Macr\`i--Toda \cite{BayerMacriToda2014}. For the Gepner charge, Toda
already proposed the first slope tilt at $-1/2$ and the second tilt using
the imaginary part of that charge
\cite[Sections~3.4--3.6 and Conjecture~3.9]{TodaQuintic2013}.
His Remark~3.10 identifies the existence of Harder--Narasimhan filtrations
for this irrational charge as a remaining difficulty.
The existence of numerical Bridgeland stability conditions on every smooth
quintic was established by Li \cite[Theorem~1.3]{Li2019}.
Theorem~\ref{thm:main} adds the specified Gepner symmetry.
Our slope-sheaf input is a consequence
of Xu's stronger Bogomolov--Gieseker inequality
\cite[Theorem~1.3]{Xu2026}, whose hypotheses and precise comparison with the
bound used here are recorded in Lemma~\ref{lem:slope-input}.

\paragraph{Main idea and proof roadmap.}
The central step is to retain two independent perturbations of Toda's second
cut. They give inequalities controlling rank and first Chern character; the
complex charge then controls the other two numerical coordinates. We obtain
this support estimate before defining semistable objects. It makes finite
refinement possible even though the image of the numerical lattice under the
charge need not be discrete.

The proof proceeds through four stages.
\begin{enumerate}
\item In Section~\ref{sec:numerical}, we record periodicity and compute the
unique normalized eigencharge. We also identify the full numerical group
with a rank-four lattice. The resulting coordinate estimate will turn
control of rank, first Chern character, and charge into full support.

\item In Section~\ref{sec:aisles}, we construct a family of bounded aisles
$U_\eta$, where an aisle is the nonpositive part of a bounded
$t$-structure and $\eta=(\eta_0,\eta_1)$ ranges over a fixed open square in
$\R^2$. The construction tilts $\Coh(X)$ twice, beginning at slope
$-1/2$, and varies the rank and first-Chern coefficients of the second cut.
Xu's inequality supplies strict margins for the comparison. The crucial
conclusion of Proposition~\ref{prop:uniform-aisle} is
\[
 \Phi U_{\eta}\subset U_{\theta}.
\]
Both parameters vary independently throughout the same square.

\item In Section~\ref{sec:grid}, periodicity and positivity refine the
translated aisles to a nested grid $V_k$, indexed by integers, with phase
spacing $1/5$ (Proposition~\ref{prop:fine-grid}). The category
$\mathcal S_k=V_k\cap V_{k+1}^\perp$ between consecutive aisles is called
a block; here $\perp$ is right $\Hom$-orthogonality. Its charges lie in a
closed sector of angle $\pi/5$. Each block object belongs to every member
of a suitable perturbed family of hearts. Testing all those perturbations
gives the full support estimate of Proposition~\ref{prop:block-support}.

\item We then refine each block into semistable factors
(Lemma~\ref{lem:block-hn}). Support and a positive projection of the charge
bound the numerical classes of subobjects and the length of any refinement.
The full lattice is discrete, so maximal phases exist and refinement stops.
We merge shared endpoint phases to obtain the slicing, construct its phase
cuts, and prove local finiteness through the exact structure of short phase
intervals (Proposition~\ref{prop:local-finiteness}).
\end{enumerate}

\paragraph{Conventions and prerequisites.}
Our aisle convention is degree $\le0$, so aisles are closed under $[1]$.
For a subcategory $\mathcal C$, the notation
$\mathcal C^\perp$ means its right $\Hom$-orthogonal. Extension closures always
include the zero object and finite iterated extensions. We use standard
coherent-sheaf cohomology, slope Harder--Narasimhan filtrations, derived
truncations, Riemann--Roch, Serre duality, and the hyperplane theorem. The
specific slope inequality is cited as above; the construction of the
stability condition from it is given in full.
